\documentclass[11pt]{article}

\usepackage[review]{acl}
\usepackage[utf8]{inputenc}
\usepackage[T1]{fontenc}
\usepackage{url}
\usepackage{booktabs}
\usepackage{enumitem}
\usepackage{amsfonts}
\usepackage{microtype}
\usepackage{xcolor}
\usepackage{xspace}
\usepackage{listings}

\title{Intent-Compiled Verified Execution for Reliable Stateful Tool Use}
\author{}
\date{}

\newcommand{\icve}{ICVE\xspace}
\lstset{
  basicstyle=\ttfamily\footnotesize,
  breaklines=true,
  columns=fullflexible,
  keepspaces=true,
  frame=single,
  xleftmargin=0pt,
  xrightmargin=0pt
}

\begin{document}
\maketitle

\begin{abstract}
Language agents increasingly operate tools that mutate persistent state, but free-form
tool-use loops often conflate intent interpretation, evidence acquisition, argument
grounding, precondition repair, action execution, and completion detection. This creates
invalid tool calls, unsafe state changes, and costly repair loops. We propose
\textbf{Intent-Compiled Verified Execution (ICVE)}, a runtime approach that compiles
state-sensitive intents into typed state machines while allowing uncovered low-risk
requests to remain ordinary model-proposed tool calls. ICVE acquires evidence
deterministically, normalizes arguments, repairs preconditions, abstains under
insufficient information, executes state transitions, and stops once postconditions are
satisfied. On public ToolSandbox no-distraction single-turn and insufficient-information
stateful tasks, ICVE eliminates invalid tool calls and unsafe state changes while using
substantially fewer model calls and tokens than ReAct. The pattern holds across weak,
main, and stronger local Qwen2.5 model scales, a local non-Qwen Phi-3-mini check, and
hosted DeepSeek replications. Component ablations show that the typed intent compiler
and abstention verifier are both necessary for the observed gains. A small real-LLM
MiniStore study shows the same direction relative to ReAct and proof-carrying tool-use,
and AppWorld shows the same verified runtime transfers: local Qwen2.5-3B, hosted
DeepSeek-chat, and hosted DeepSeek-reasoner intent extraction preserve 72/72 success on
a targeted stateful slice, while deterministic ICVE and DeepSeek-chat intent extraction
also reach 57/57 on the local official AppWorld 0.2.0 dev split with the packaged
evaluator, compared with 45/57 for the official DeepSeek ReAct-code agent. On the
local AppWorld 0.2.0 \texttt{test\_normal.txt} file, deterministic ICVE and
DeepSeek-chat intent extraction reach 164/167 overall and 164/164 supported success,
with the three failures limited to guarded abstentions.
\end{abstract}

\section{Introduction}

Stateful tool use differs from stateless API answering. A single wrong argument can delete
the wrong contact, send a message under the wrong precondition, or repeat an already
completed state transition until the execution environment raises an error. Existing
tool-use agents often represent this entire process as a text loop: the language model
chooses a tool, fills arguments, interprets observations, repairs errors, and decides when
to stop. ReAct-style agents, tool-learning systems, and recent interactive benchmarks
make this setting increasingly realistic \citep{yao2022react,schick2023toolformer,
lu2024toolsandbox,appworld2024,wang2025dialogtool}, but they leave open where stateful
execution should be verified.

Our central claim is that many failures are not reasoning failures alone. They are
\emph{execution-boundary} failures. The model is asked to perform intent compilation,
state tracking, evidence lookup, precondition management, and completion detection through
plain text. Increasing model scale helps but does not remove invalid calls or unsafe
state changes in insufficient-information settings.

ICVE changes the execution boundary. The LLM may still propose actions for uncovered
requests, but risky state-sensitive intents are handled by a runtime compiler that emits
verified tool traces. The runtime owns typed slots, evidence acquisition, safe
precondition repair, abstention, state transitions, and postcondition-based stopping.
The unit of reuse is therefore not a prompt template but a small executable contract:
once a stateful task family has a machine, weaker or stronger models only choose a frame,
while the same runtime executes and checks the state transition. This paper constructs
the machines manually; the claim is an explicit, auditable, reusable boundary once
registered, not solved machine synthesis.

The distinction is not merely stricter schema validation. A function schema can reject an
ill-typed call, but it does not decide which evidence must be acquired before a mutation,
whether an earlier transition already satisfied the goal, or whether missing state should
force abstention. Nor is ICVE only semantic parsing from utterances to API calls: the
compiled object is an executable state machine whose transition relation is checked
against a live ledger. Finally, unlike proof-carrying prompting, the runtime does not ask
the model to author the proof obligation after proposing an action. The proof burden is
compiled into the machine and reused across turns and models.

This paper makes three contributions:
\begin{enumerate}[leftmargin=*]
\item A typed intent/state-machine abstraction for stateful tool use. Each risky task
family is represented as an \texttt{IntentMachine} with an explicit schema, compiler, and
handler.
\item A verified runtime that owns evidence acquisition, argument grounding, precondition
repair, insufficient-information abstention, and completion detection.
\item A public-benchmark evaluation, including ToolSandbox, a targeted AppWorld transfer
slice, and a same-split local official AppWorld dev comparison, showing that
runtime-owned execution sharply reduces invalid calls and unsafe state changes while
avoiding the repair overhead that made proof-carrying tool-use unattractive in the same
setting.
\end{enumerate}

\begin{figure}[t]
\centering
\scriptsize
\setlength{\fboxsep}{0.7pt}
\newcommand{\boundbox}[3]{\fcolorbox{black!45}{#1}{\parbox[c][0.34cm][c]{#2}{\centering #3}}}
\begin{tabular}{@{}p{0.44\columnwidth}@{\hspace{0.012\columnwidth}}p{0.44\columnwidth}@{}}
\multicolumn{1}{c}{\textbf{Free-form loop}} &
\multicolumn{1}{c}{\textbf{\icve state machine}}\\[-0.15em]
\fcolorbox{red!40}{red!4}{\parbox[t][2.23cm][t]{0.415\columnwidth}{
\centering
\textit{model owns search, mutate, stop}\\[-0.1em]
\boundbox{red!9}{0.33\columnwidth}{guess evidence / args}\\[-0.05em]
\(\downarrow\)\\[-0.45em]
\boundbox{red!9}{0.33\columnwidth}{mutating tool call}\\[-0.05em]
\(\downarrow\)\\[-0.45em]
\boundbox{red!9}{0.33\columnwidth}{parse error and retry}\\[-0.05em]
\(\leftarrow\) unsafe loop
}} &
\fcolorbox{blue!40}{blue!3}{\parbox[t][2.23cm][t]{0.415\columnwidth}{
\centering
\textit{model-owned}\\[-0.1em]
\boundbox{blue!8}{0.33\columnwidth}{request \(\to\) IntentFrame}\\[-0.05em]
\(\Downarrow\) \textit{boundary shift}\\[-0.35em]
\textit{runtime-owned}\\[-0.1em]
\boundbox{green!9}{0.33\columnwidth}{schema + ledger}\\[-0.05em]
\(\downarrow\)\\[-0.45em]
\boundbox{green!9}{0.33\columnwidth}{repair or abstain}\\[-0.05em]
\(\downarrow\)\\[-0.45em]
\boundbox{green!9}{0.33\columnwidth}{action + stop}
}}
\end{tabular}
\caption{The execution-boundary shift. A free-form agent owns evidence search, mutation,
error recovery, and stopping, so missing state can become invalid calls or unsafe loops.
\icve restricts the model to typed frame selection and moves ledger checks, repair,
abstention, tool execution, and postcondition stopping into the runtime.}
\label{fig:icve-overview}
\end{figure}

\section{Method}

\subsection{Typed Intent Machines}

ICVE represents each covered stateful task family as an intent machine:
\begin{lstlisting}
Machine  = (schema, compiler, handler)
schema   = (intent_type, slots, terminal_states)
compiler = request x dialogue x tools -> frame or None
handler  = frame x ledger x tools -> tool_action or final_action
\end{lstlisting}

An \texttt{IntentFrame} is the runtime-owned state for a user intent:
\begin{lstlisting}
IntentFrame = (
  intent_type,
  slots: name -> (value, required, source),
  state,
  missing_slots,
  abstain_reason
)
\end{lstlisting}

The compiler is deliberately typed and conservative. It does not ask the model to emit a
proof or a full tool trace. Instead, it maps a user request into a small frame whose slot
values are extracted from the request, copied from prior observations, or produced by a
canonicalizer. In the ToolSandbox binding, the registry currently covers 13 intent
families: currency conversion, stock lookup, reverse geocode, current city, last-message
contact update, bulk contact relationship update, contact name lookup, location phone
lookup, holiday timestamp/countdown, weather, distance, and reminder.

The machine interface is also the method's threat boundary. A model may choose the wrong
intent type or fill a wrong slot; that risk is handled by schema validation, conservative
compilers, and domain verifiers. The model is not allowed to choose arbitrary
state-changing calls for a covered intent. For those intents, handlers may only emit
runtime actions derived from validated slots and ledger observations, and a policy
verifier can reject or rewrite the action before execution. Thus ICVE's claim is
conditional and inspectable: correctness depends on the registered machine and verifier,
not on an opaque promise that the model will reason safely.

\subsection{Compilation}

Compilation proceeds by registry order. Figure~\ref{fig:icve-overview} illustrates the
resulting boundary: the compiler selects a frame, while the runtime owns state checks,
repair, abstention, execution, and stopping.
On each turn, ICVE first runs an insufficient-information verifier. If required tools or
information are unavailable, the verifier returns an abstention frame. Otherwise, the
runtime tries registered compilers until one returns a frame, and validates that frame
against its schema. If no registered machine applies, control falls back only for
uncovered low-risk requests.

This loop is implemented in a benchmark-agnostic runtime named \texttt{RaveRuntime} in
the artifact for historical reasons. Domain bindings provide intent machines plus a small
runtime policy: a state-ledger interface
for grounded observations and a verifier interface for final-message wording and
transition rejection. The generic runtime owns the registry, uniqueness checks, frame
validation, completion/abstention hooks, and handler dispatch. A policy verifier can
rewrite a proposed \texttt{ToolAction} into a \texttt{FinalAction}, which gives each
domain binding a single runtime boundary for rejecting unsafe or previously failed
transitions.

\subsection{Verified State-Machine Runtime}

The runtime consumes a state ledger rather than trusting raw dialogue text. The ledger
records user-provided values, successful tool observations, and recent tool errors; a
ToolSandbox ledger instantiates this interface with settings, contacts, messages, and
reminders. A state-machine handler consumes the current \texttt{IntentFrame} and ledger,
then emits exactly one transition:
\begin{verbatim}
ToolAction  = (tool, args, reason)
FinalAction = (message, reason)
\end{verbatim}

For read-only actions, the verifier normalizes arguments and blocks empty or ungrounded
searches. For state-changing actions, it checks schema conformance, argument grounding,
state preconditions, and duplicate completion before allowing exactly one transition. If
the requested postcondition is already satisfied, the runtime stops. When a safe repair
exists, such as disabling low-battery mode before enabling Wi-Fi, the runtime emits the
repair action before the requested mutation. When repair would require missing
information or unavailable tools, it abstains instead of letting the model guess.
Thus the verifier is stricter than a call schema: it checks the transition's evidence,
preconditions, and completion state, not only the argument types of one invocation.

The design point is that stateful execution is runtime-owned. The model can still be used
for uncovered or low-risk requests, but registered stateful intents are compiled and
executed through deterministic state machines.

\section{Experimental Setup}

We evaluate under real local LLM inference rather than oracle policy traces. The main
agent LLM is served through a local OpenAI-compatible CUDA server. ToolSandbox tools that
would otherwise require external RapidAPI calls use deterministic offline fixtures so the
benchmark remains repeatable.

All reported local-model experiments use temperature 0 inference on a single NVIDIA RTX
3080 Ti GPU with 12GB memory. Qwen2.5-7B and Phi-3-mini are run with 4-bit quantization;
Qwen2.5-0.5B and Qwen2.5-3B are run as local instruction models through the same
OpenAI-compatible interface. No model training or fine-tuning is performed.

For the hosted replication, we run the same ToolSandbox suites with the DeepSeek API
model identifiers \texttt{deepseek-chat} and \texttt{deepseek-reasoner} through the same
OpenAI-compatible client. API keys are loaded from environment variables and are not
written to command lines or result files.

For AppWorld, we use two settings. First, we run an isolated \texttt{pctu-appworld}
environment on a targeted public stateful slice covering 72 tasks across twenty-four
registered AppWorld intent machines; some hosted DeepSeek rows are combined from an
earlier 66-task run plus the final six-task Spotify-family expansion. Second, we run
ICVE on the complete local official AppWorld 0.2.0 \texttt{dev.txt} split in a separate
\texttt{pctu-appworld-agents} environment and evaluate with AppWorld's packaged
evaluator. The deterministic path uses typed compilers only; the real-LLM paths use
local Qwen2.5-3B, DeepSeek-chat, or DeepSeek-reasoner to produce a typed intent frame
and then execute that frame through the same verified runtime. AppWorld handlers
interact only through the live \texttt{apis} object; the runtime policy blocks code that
references task files, ground truth, compiled solutions, or oracle data.

Metrics are task success, ToolSandbox similarity, invalid tool calls per task, unsafe
state changes per task, verifier rejections, runtime repair actions, model calls, and a
token proxy computed from local model usage. The kill criterion is not only safety: a
candidate must reduce invalid calls and unsafe state changes without letting repair cost
erase task-success or token-efficiency gains.
For AppWorld, we report evaluator success, unsafe state-test failures, failed API repair
attempts, API calls, model calls, and model-reported token usage.

We use Wilson 95\% intervals over the fixed episode counts in each main suite; the
complete interval table is archived in the supplement. Invalid and unsafe counts are
reported as per-episode means on the same fixed slices.

The experiments answer three questions: whether ICVE reduces invalid or unsafe state
changes, whether the same boundary transfers across models and environments, and whether
typed machines plus abstention are necessary beyond intent recognition.
The anonymous supplement includes source code, task-id lists, evaluator summaries, and
commands for reproducing the reported rows.

\paragraph{Baselines.}
We compare against ReAct \citep{yao2022react}, a plain JSON tool-use loop with no runtime
intent compiler; PCTU, a proof-carrying ablation that asks the model to attach action
preconditions and lets a verifier reject unsafe steps; and ICVE ablations that remove the
typed DSL or the abstention verifier.

\section{Results}

\subsection{MiniStore Diagnostic}

MiniStore is a small local stateful environment used as a diagnostic rather than as the
headline benchmark. It shows why model-authored proof obligations are not enough.

\begin{table}[h]
\centering
\scriptsize
\setlength{\tabcolsep}{3pt}
\begin{tabular}{lrrrrrr}
\hline
method & succ. & goal & invalid & unsafe & LLM & tokens \\
\hline
ReAct & 0.375 & 1.000 & 1.125 & 1.125 & 5.500 & 2477.750 \\
PCTU & 0.250 & 0.250 & 0.000 & 0.000 & 7.250 & 3169.375 \\
ICVE & 1.000 & 1.000 & 0.000 & 0.000 & 2.250 & 603.875 \\
\hline
\end{tabular}
\caption{MiniStore real-LLM diagnostic with Qwen2.5-3B. PCTU suppresses invalid and
unsafe actions but pays too much proof/repair overhead.}
\label{tab:ministore}
\end{table}

\subsection{Public ToolSandbox}

The strongest evidence is on ToolSandbox \citep{lu2024toolsandbox}. Table
\ref{tab:toolsandbox-main} reports the primary Qwen2.5-3B results, and
Figure~\ref{fig:toolsandbox-deep} summarizes the scale and failure-mode effects.

\begin{figure}[t]
\centering
\scriptsize
\setlength{\tabcolsep}{2pt}
\newcommand{\ratebar}[2]{\fcolorbox{black!15}{#1}{\rule{#2}{0.12cm}}}
\newcommand{\stackseg}[2]{{\color{#1}\rule{#2}{0.15cm}}}
\begin{tabular}{@{}lccc@{}}
\multicolumn{4}{@{}l}{\textbf{(a) Success across model scales}}\\[-0.15em]
\toprule
model & ReAct & \icve & closed gap \\
\midrule
0.5B & \ratebar{red!22}{0.05cm} .03 & \ratebar{blue!28}{1.35cm} 1.00 & +.97 \\
3B & \ratebar{red!22}{0.10cm} .07 & \ratebar{blue!28}{1.35cm} 1.00 & +.93 \\
7B & \ratebar{red!22}{0.20cm} .13 & \ratebar{blue!28}{1.35cm} 1.00 & +.87 \\
DeepSeek & \ratebar{red!22}{0.18cm} .12 & \ratebar{blue!28}{1.35cm} 1.00 & +.88 \\
\bottomrule
\end{tabular}
\vspace{0.25em}

\begin{tabular}{@{}lcl@{}}
\multicolumn{3}{@{}l}{\textbf{(b) Insuff.-info terminations}}\\[-0.15em]
\toprule
method & distribution & mix \\
\midrule
ReAct & \stackseg{blue!65}{1.15cm}\stackseg{red!70}{0.99cm}\stackseg{gray!55}{0.16cm}
& .50 / .43 / .07 \\
\icve & \stackseg{blue!65}{2.30cm}
& 1.00 / .00 / .00 \\
\bottomrule
\end{tabular}
\caption{Deep ToolSandbox analysis. ICVE closes the success gap across weak, main, and
hosted model scales and changes insufficient-information failures by eliminating unsafe
state transitions. Blue denotes success, red unsafe mutation, and gray other invalid or
incomplete termination.}
\label{fig:toolsandbox-deep}
\end{figure}

\begin{table*}[t]
\centering
\scriptsize
\setlength{\tabcolsep}{4pt}
\begin{tabular}{llrrrrrrr}
\toprule
suite & method & n & succ. & invalid & unsafe & repair & LLM & tokens \\
\midrule
single-turn & ReAct & 73 & 0.0685 & 2.1507 & 0.0000 & 0.0000 & 11.8219 & 11048.0 \\
single-turn & ICVE no DSL & 73 & 0.4932 & 0.3288 & 0.0000 & 0.8082 & 3.9726 & 3307.9 \\
single-turn & \icve & 73 & 1.0000 & 0.0000 & 0.0000 & 2.4658 & 0.1233 & 32.7 \\
insufficient-info & ReAct & 28 & 0.5000 & 2.1786 & 0.4286 & 0.0000 & 11.7500 & 9896.4 \\
insufficient-info & ICVE no abstention & 28 & 0.5000 & 1.5714 & 0.4286 & 0.3929 & 7.2857 & 6612.1 \\
insufficient-info & \icve & 28 & 1.0000 & 0.0000 & 0.0000 & 1.5714 & 0.0000 & 0.0 \\
\bottomrule
\end{tabular}
\caption{Main public ToolSandbox results with Qwen2.5-3B. Metrics after success are
per-episode means. ICVE performs deterministic repair actions but still uses far fewer
model calls and tokens; the typed DSL is necessary on single-turn tasks, and the
abstention verifier is necessary on insufficient-information tasks.}
\label{tab:toolsandbox-main}
\end{table*}

The Wilson 95\% confidence intervals for the primary success rates are separated from the
mean-count columns: on the 73-episode single-turn suite, ICVE is 1.0000
[$0.9500, 1.0000$] while ReAct is 0.0685 [$0.0296, 0.1505$]; on the 28-episode
insufficient-information suite, ICVE is 1.0000 [$0.8794, 1.0000$] while ReAct and the
no-abstention ablation are both 0.5000 [$0.3263, 0.6737$].

\subsection{Model Scale}

The effect is not a Qwen2.5-3B-only artifact. Table \ref{tab:model-scale} summarizes weak
and stronger local model runs. The 7B result uses a local 4-bit quantized model and should
not be interpreted as frontier-model coverage.

\begin{table*}[t]
\centering
\scriptsize
\setlength{\tabcolsep}{4pt}
\begin{tabular}{lllrrrrr}
\toprule
model & suite & method & n & succ. & invalid & unsafe & tokens \\
\midrule
0.5B & single-turn & ReAct & 30 & 0.0333 & 5.6667 & 0.0000 & 8054.5 \\
0.5B & single-turn & \icve & 30 & 1.0000 & 0.0000 & 0.0000 & 54.3 \\
0.5B & insufficient-info & ReAct & 28 & 0.6071 & 3.2500 & 0.3214 & 8367.3 \\
0.5B & insufficient-info & \icve & 28 & 1.0000 & 0.0000 & 0.0000 & 0.0 \\
7B-4bit & single-turn & ReAct & 30 & 0.1333 & 0.4333 & 0.0000 & 3483.5 \\
7B-4bit & single-turn & \icve & 30 & 1.0000 & 0.0000 & 0.0000 & 34.6 \\
7B-4bit & insufficient-info & ReAct & 28 & 0.4286 & 1.2143 & 0.3929 & 4894.4 \\
7B-4bit & insufficient-info & \icve & 28 & 1.0000 & 0.0000 & 0.0000 & 0.0 \\
\bottomrule
\end{tabular}
\caption{ICVE preserves the zero-invalid, zero-unsafe pattern across weak and stronger
local Qwen2.5 scales; the primary Qwen2.5-3B results appear in Table
\ref{tab:toolsandbox-main}.}
\label{tab:model-scale}
\end{table*}

\subsection{Cross-Family Diagnostic}

We also ran a small non-Qwen diagnostic with Phi-3-mini. Table
\ref{tab:phi3-diagnostic} reports 10 core ToolSandbox scenarios and 10
insufficient-information scenarios. This is not frontier-model evidence, but it addresses
the narrow concern that the result is only a Qwen-family artifact.

\begin{table*}[t]
\centering
\scriptsize
\setlength{\tabcolsep}{4pt}
\begin{tabular}{llrrrrr}
\toprule
suite & method & n & succ. & invalid & unsafe & tokens \\
\midrule
core10 & ReAct & 10 & 0.0000 & 1.8000 & 0.0000 & 6072.0 \\
core10 & \icve & 10 & 1.0000 & 0.0000 & 0.0000 & 87.9 \\
insufficient10 & ReAct & 10 & 0.6000 & 1.7000 & 0.3000 & 5216.3 \\
insufficient10 & ICVE no abstention & 10 & 0.6000 & 0.9000 & 0.4000 & 1123.2 \\
insufficient10 & \icve & 10 & 1.0000 & 0.0000 & 0.0000 & 0.0 \\
\bottomrule
\end{tabular}
\caption{Phi-3-mini cross-family diagnostic. The typed runtime remains clean on both
slices, and removing the abstention verifier again permits invalid and unsafe attempts.}
\label{tab:phi3-diagnostic}
\end{table*}

\subsection{Hosted OpenAI-Compatible Replication}

We also replicate the primary ToolSandbox suites with the hosted DeepSeek-chat endpoint.
Table \ref{tab:deepseek-replication} shows the same qualitative pattern: ICVE reaches
1.0000 success with zero invalid calls and zero unsafe state changes on both suites, while
ReAct remains substantially less successful and more expensive. The no-DSL and
no-abstention ablations again expose the two mechanisms: typed compilation is needed on
single-turn stateful tasks, and abstention verification is needed when the user omits
required information.

\begin{table*}[t]
\centering
\scriptsize
\setlength{\tabcolsep}{4pt}
\begin{tabular}{llrrrrr}
\toprule
suite & method & n & succ. & invalid & unsafe & tokens \\
\midrule
single-turn & ReAct & 73 & 0.1233 & 0.5342 & 0.0000 & 4397.1 \\
single-turn & ICVE no DSL & 73 & 0.5479 & 0.0548 & 0.0000 & 1258.6 \\
single-turn & \icve & 73 & 1.0000 & 0.0000 & 0.0000 & 13.6 \\
insufficient-info & ReAct & 28 & 0.3214 & 1.2143 & 0.6071 & 7299.6 \\
insufficient-info & ICVE no abstention & 28 & 0.5000 & 0.6786 & 0.4286 & 3234.4 \\
insufficient-info & \icve & 28 & 1.0000 & 0.0000 & 0.0000 & 0.0 \\
\bottomrule
\end{tabular}
\caption{Hosted DeepSeek-chat ToolSandbox replication. Metrics after success are
per-episode means; full outputs are archived in the artifact manifest.}
\label{tab:deepseek-replication}
\end{table*}

We also ran DeepSeek-reasoner on the same two suites; both full runs passed the kill
criteria, and the recorded summaries are listed in the artifact manifest.

\subsection{AppWorld}

The same verified runtime transfers to AppWorld on a targeted public state-changing
slice. The active rows cover 72 tasks across twenty-four registered AppWorld
intent machines spanning phone, Venmo, Gmail, Simple Note, Spotify, Amazon, alarms,
Bucket List, and cross-app sharing tasks.

\begin{table*}[t]
\centering
\scriptsize
\setlength{\tabcolsep}{4pt}
\begin{tabular}{lrrrrrr}
\toprule
method & n & succ. & inv./repair & unsafe & LLM & tokens \\
\midrule
deterministic compiler & 72 & 72/72 & 0.0694 & 0.0000 & 0.0000 & 180.5 \\
Qwen2.5-3B intent & 72 & 72/72 & 0.0694 & 0.0000 & 1.0000 & 3220.3 \\
DeepSeek-chat intent & 72 & 72/72 & 0.0694 & 0.0000 & 1.0000 & 3352.2 \\
DeepSeek-reasoner intent & 72 & 72/72 & 0.0694 & 0.0000 & 1.0000 & 3462.3 \\
Qwen direct code & 72 & 0/72 & 1.0000 & 0.0000 & 1.0000 & 1544.0 \\
Qwen typed intent/code & 72 & 0/72 & 1.0000 & 0.0000 & 2.0000 & 4674.8 \\
Qwen code repair & 72 & 0/72 & 2.7083 & 0.0000 & 2.8889 & 5342.8 \\
Qwen multi-step code & 72 & 0/72 & 3.7500 & 0.0000 & 4.3472 & 9031.7 \\
DeepSeek direct code & 72 & 4/72 & 0.8611 & 0.0139 & 1.0000 & 2028.9 \\
DeepSeek typed intent/code & 72 & 7/72 & 0.7917 & 0.0139 & 2.0000 & 5478.9 \\
DeepSeek code repair & 72 & 21/72 & 1.5556 & 0.0417 & 2.3472 & 5898.8 \\
DeepSeek multi-step code & 72 & 20/72 & 1.2222 & 0.0000 & 4.3056 & 12128.6 \\
\bottomrule
\end{tabular}
\caption{Targeted public AppWorld slice. Metrics after success are per-task means. All
rows use the same 72-task dataset; hosted DeepSeek rows are combined from incrementally
flushed shards where needed. AppWorld's packaged evaluator confirms the corresponding
task-completion counts: 72/72 for ICVE runtime rows, 0/72 for local Qwen code baselines,
4/72 for DeepSeek direct code, 7/72 for DeepSeek typed intent/code, 21/72 for DeepSeek
repair, and 20/72 for DeepSeek multi-step code.}
\label{tab:appworld-slice}
\end{table*}

The AppWorld result separates intent recognition from verified execution. Typed-intent
code gives the model the same frame but still asks it to write the mutation code; it
reaches 0/72 with Qwen2.5-3B and 7/72 with DeepSeek-chat. Direct code and multi-step
code-observation variants likewise fail or become expensive, while DeepSeek repair reaches
21/72 but introduces unsafe changes. The error pattern is consistent: models often know
the requested operation but fail at grounded execution. One unverified DeepSeek-chat
intent run extracted \texttt{archive} for a delete instruction; the instruction-aware
slot verifier rewrites the frame before execution. ICVE's residual repair cost is
concentrated in Venmo transfers, where public metadata reveals card expiry but not hidden
balances, so the runtime probes non-expired cards and then stops after success.

We then moved the ICVE binding to the complete local official AppWorld 0.2.0
\texttt{dev.txt} split. This adds the remaining official dev task families as registered
intent machines, including navigation-based Spotify artist lookup, Venmo social-feed
aggregation, file-prefix moves, current-artist follower lookup, and Simple Note markdown
export. Table~\ref{tab:appworld-dev57} reports AppWorld's packaged evaluator on all 57
dev tasks. Both deterministic \icve and \icve with DeepSeek-chat intent extraction reach
57/57 local success and 100.0 task-goal / 100.0 scenario-goal completion, while the
official DeepSeek ReAct-code agent reaches 79.0 / 73.7 and 45/57. The residual
0.0351 invalid API attempts per task are two Venmo payment-card probes: public card
metadata exposes expiry but not balance, so the runtime tries one non-expired
insufficient-balance card and then succeeds with the next public card. We do not use the
private admin balance API to remove this residual repair cost.

\begin{table*}[t]
\centering
\scriptsize
\setlength{\tabcolsep}{4pt}
\begin{tabular}{lrrrrrr}
\toprule
method & n & task & scen. & invalid & unsafe & LLM \\
\midrule
deterministic compiler & 57 & 100.0 & 100.0 & 0.0351 & 0.0000 & 0.0000 \\
DeepSeek-chat intent & 57 & 100.0 & 100.0 & 0.0351 & 0.0000 & 1.0000 \\
official DeepSeek ReAct-code & 57 & 79.0 & 73.7 & -- & -- & 15.4211 \\
\bottomrule
\end{tabular}
\caption{Complete local official AppWorld 0.2.0 dev split for ICVE rows. Task and
scenario columns are packaged-evaluator percentages; invalid, unsafe, and LLM
are per-task means from the local runtime logs for ICVE and official runner logs for
ReAct-code. The official ReAct-code row uses 6.45M tokens total.}
\label{tab:appworld-dev57}
\end{table*}

The official DeepSeek ReAct-code dev57 run completes all tasks in 75.9 minutes, with
879 LLM calls and 6.45M tokens. For context, its earlier \texttt{dev10\_full50} slice
reached 90.0 / 75.0 with 174 model calls and 1.40M tokens, while ICVE rows reached
100.0/100.0. These are same-split local official-dev comparisons, not held-out test or
leaderboard submissions.

As a held-out smoke, we then expanded beyond the official dev split to the beginning of
AppWorld \texttt{test\_normal}. The first-12 smoke initially exposed zero coverage; after
registering four typed machines, deterministic ICVE reached 12/12 with zero invalid
tool calls and zero unsafe state changes. We further expanded the smoke to the first 120
\texttt{test\_normal} tasks. After registering additional Simple Note and Venmo
machines plus Simple Note-to-Spotify playlist, release-month, reunion-RSVP file,
phone-sourced Spotify playlist update, roommate dinner-settlement, offline playback, and
Venmo card-path repair machines, the registry supports 117/120 and succeeds on all
117/117 supported tasks, giving 117/120 overall success, 0 unsafe state changes, and
0.0250 invalid/tool per task. The invalid count consists of 3 unsupported or no-action
abstentions from the unsupported vacation-settlement family. The Venmo named-payment
family now adds an available public supervisor card before transfer instead of probing
failing cards, and the roommate payment-request family rotates successful non-expired
cards to avoid reusing a depleted card first. We report this as held-out-prefix smoke
evidence rather than a leaderboard result. DeepSeek-chat intent extraction on the same
first 120 tasks matches this deterministic coverage profile after the instruction-aware
abstention verifier is applied: 117/120 overall success, 117/120 supported, 117/117
supported success, 0 unsafe
state changes, 0.0250 invalid/tool per task, and 1 LLM call per task. A strict verifier
blocks unsupported vacation-settlement tasks before code execution and counts them as
invalid no-action results. Earlier LLM
over-generalizations on unsupported housing-bill, offline-playback, and carpool Venmo
tasks are now either converted to abstention or replaced by dedicated typed machines
rather than similar-looking state changes. Running the remaining local
\texttt{test\_normal.txt} entries gives the same profile at file level: both deterministic
\icve and DeepSeek-chat intent extraction reach 164/167 overall, 164/167 supported, and
164/164 supported success with 0 unsafe state changes; the only failures are the same
three protected vacation-settlement abstentions. The 121--167 continuation covers the
remaining 47 entries in the local file and reaches 47/47 overall success and 47/47
supported success, with 0 unsafe state changes and 0 invalid tool calls. A full
DeepSeek-chat intent-extraction companion run on the same 47 tasks matches this result
with 1 LLM call per task, 0 unsafe state changes, and 0 invalid tool calls. This
file-level continuation remains a coverage diagnostic rather than a full held-out test
split or leaderboard result.
We also ran a separate beginning-of-\texttt{test\_challenge} Amazon diagnostic after
adding two conservative saved-list machines: the first 24 tasks improved from 3/24 to
12/24 success with zero unsafe state changes, while product-search purchase requests
that require choosing among sellers, dimensions, or multiple candidate products remain
unsupported.

\subsection{Multi-Turn Diagnostic}

On a 28-episode hidden-task multi-turn diagnostic with Qwen2.5-0.5B, ICVE reaches
0.8929 success, 0.9773 mean similarity, 0 invalid calls, and 0 unsafe state changes.
This is an extension result, not the headline claim: a small number of recency and
benchmark-interface near-misses remain.

\section{Analysis}

The empirical pattern is that free-form agents overpay for tasks with simple state
semantics. ReAct must repeatedly infer which evidence is needed, call tools with fragile
arguments, interpret observations, decide whether preconditions are satisfied, and detect
completion. PCTU improves safety by rejecting bad calls, but it still relies on the model
to propose proof-bearing actions. ICVE moves these repeated proof obligations into the
runtime. Runtime actions increase, but model calls and token usage collapse because the
agent avoids repeated model-mediated repair.

\section{Related Work}

ReAct-style agents interleave reasoning and action through text \citep{yao2022react}.
Reflexion adds verbal self-feedback \citep{shinn2023reflexion}, and Toolformer explores
learning tool use from self-supervised signals \citep{schick2023toolformer}. Recent
tool-learning systems and datasets further scale API selection and function calling
\citep{li2023apibank,shen2024smalltoollearners,qin2024toolllm,patil2024gorilla,
liu2025toolace,dang2025guidedtemplates}. These lines improve model behavior, but they
still leave stateful execution largely model-owned. ICVE instead changes the execution
boundary for covered risky intents.

Tool-use benchmarks have become more realistic. ToolSandbox introduces stateful,
interactive evaluation with insufficient-information and canonicalization cases
\citep{lu2024toolsandbox}. AppWorld, BFCL, and tau-bench broaden interactive and
function-calling evaluation \citep{appworld2024,bfcl2024,taubench2024}; broader agent
benchmarks stress multi-step interaction in text, web, and operating-system environments
\citep{liu2024agentbench,zhou2024webarena,xie2024osworld}. Recent stateful-dialogue
tool-use benchmarks further show that state tracking remains difficult
\citep{wang2025dialogtool}. ICVE is complementary: it is a runtime method intended to
reduce invalid and unsafe execution in such environments.

Guardrails, function schemas, and sandboxed risk evaluation validate or probe tool-agent
behavior \citep{ruan2024toolemu,xie2025toolsafety,faghih2025toolpreferences}. ICVE
verifies traces against intent-level state machines: the runtime checks not only whether
one call is syntactically valid, but whether it is grounded in the requested state change
and whether the postcondition has already been satisfied. This shifts the unit of
verification from one tool invocation to a stateful intent. The AppWorld typed-intent-code
ablation highlights the distinction: giving the model a valid frame but letting it write
the mutation code still fails on most tasks, because the difficult part is not only
recognizing the intent but executing the grounded state transition.

Semantic parsing and task-oriented dialogue systems also use intent and slot structure,
but their usual target is interpreting user utterances or filling API arguments. ICVE
uses typed intent frames as an execution substrate: compilers produce state-machine
frames, handlers acquire missing evidence and repair preconditions, and the verifier can
abstain when the state needed for safe execution is unavailable. This is why the ablation
without the typed DSL loses both success and invalid-call control in Table
\ref{tab:toolsandbox-main}.

Proof-carrying tool-use approaches ask the model to make safety obligations explicit.
Our MiniStore diagnostic shows the limitation of that boundary: proof checks can suppress
invalid and unsafe mutations, but the model still pays for repeated proof generation,
rejection, and repair. ICVE keeps the verification boundary while moving recurring proof
obligations into deterministic machines.

\section{Conclusion}

ICVE reframes reliable stateful tool use as verified execution, not as longer prompting
or model-authored proofs. The agent supplies intent; a typed runtime owns risky state
transitions. On the covered public ToolSandbox slices, this boundary reduces invalid
calls and unsafe state changes while improving task success and reducing repair cost.
AppWorld transfer, official-dev, and local file-level runs show the same execution
boundary beyond ToolSandbox. Public leaderboard evaluation and broader real-LLM coverage
remain the next empirical step.

\section*{Limitations}

The current strongest claim is intentionally narrow. The clean real-LLM evidence is on
ToolSandbox no-distraction single-turn and insufficient-information suites plus targeted
AppWorld transfer, complete local official AppWorld dev57, and local
\texttt{test\_normal.txt} file-level runs. These runs use AppWorld's packaged evaluator
where applicable, but they are not a public AppWorld leaderboard submission. Unsupported
vacation-settlement tasks remain blocked as invalid no-action results rather than being
forced through a similar-looking Venmo machine. Multi-turn coverage is promising but not
yet broad enough for a headline claim. The current IntentMachine registry is explicit and
partly hand-engineered; future work should learn or synthesize candidate machines and
then verify them.

The runtime does not eliminate the need for a language model. It changes where the model
is trusted. ICVE still needs the model for uncovered intents, natural-language
clarification, and domains whose state machines have not been registered. The safety claim
therefore applies to covered high-risk intents, not arbitrary tool execution.

Finally, hosted-model coverage remains narrow: the evidence spans Qwen2.5 scales, a
small Phi-3-mini diagnostic, and hosted DeepSeek replications, but not a broad set of
proprietary frontier models.

\section*{Ethical Considerations}

Reliable stateful tool execution has positive applications in reducing accidental
mutations, repeated tool calls, and unsafe guessing by deployed agents. The same runtime
discipline could also make high-impact automated systems more capable, so the method
should be paired with domain authorization, audit logs, and tool-level access controls.
ICVE is not a substitute for policy enforcement: it verifies covered state transitions
against declared machine-readable rules, but cannot guarantee safety for tools or intents
outside the registry.

\section*{Reproducibility Statement}

All reported benchmark results are produced by local scripts and archived output
directories listed in the artifact manifest. For AppWorld, the implementation uses public
task instructions, public API documentation, and live public app state observed through
the benchmark APIs. It does not read hidden answers, private evaluation code, oracle
solutions, or compiled solutions when constructing actions. Unsupported or ambiguous
high-risk families are reported as abstentions rather than forced executions. The
anonymous supplement includes source code, task-id lists, evaluator summaries, and
environment notes, but excludes API keys and hosted-model credentials. It also avoids
redistributing AppWorld's protected bundle, raw databases, access tokens, raw task logs,
or downloaded data bundles; users should obtain AppWorld through its official release
channel and follow its license and redistribution terms.

\bibliography{rave_refs}

\end{document}
